\subsection{章末確認問題}
1. プロセス、活動、タスクの違いを説明せよ。

2. この章でいう「プロセス」が、プログラム実行時の処理ではない理由を説明せよ。

3. ライフサイクルが、単なるプロセス一覧よりも有用である理由を説明せよ。

4. 技術プロセス、技術管理プロセス、組織的プロジェクト支援プロセス、合意プロセスの違いを説明せよ。

5. 予測型、反復型、進化型、増分型、継続的開発を比較せよ。

6. アジャイルは文書不要という意味ではない。なぜそう言えるか。

7. DevOps が必要になった背景を、リリース頻度と組織連携の観点から説明せよ。

8. ライフサイクル適応で、判断と根拠を文書化する理由を説明せよ。

9. GQM の「目標、質問、メトリクス」の順序が重要である理由を説明せよ。

10. アジャイルのふりかえりが、プロセス改善の学習ループである理由を説明せよ。

\begin{pointbox}{解答の方向性}1 は入力・活動・出力と作業粒度の違い、2 は人間の作業活動を扱う点、3 はプロセス間依存と制約を含める点を説明する。5 は要求変更、見積り、成果物の出し方で比較する。9 は測れるものを集めるのではなく、改善目標から測定を設計する点を説明する。\end{pointbox}
